\documentclass[12pt]{article}
\usepackage{amsmath, amssymb}
\usepackage[margin=1in]{geometry}
\setlength{\parskip}{6pt}
\title{QUBO Formulation: Steiner Tree Problem}
\date{}
\begin{document}
\maketitle

\subsection*{Steiner Tree Problem}

\paragraph{Problem Statement.}
Given a weighted undirected graph $G = (V, E)$ with non-negative edge weights $w_e$ for each $e \in E$, and a distinguished subset of terminal vertices $T \subseteq V$, the Steiner Tree Problem seeks a minimum-weight connected subgraph of $G$ that spans all vertices in $T$. The goal is to select a subset of edges such that all terminals lie in the same connected component, while minimizing the sum of the weights of the selected edges.

\paragraph{Decision Variables.}
For each edge $e \in E$, introduce a binary decision variable $x_e \in \{0, 1\}$, where $x_e = 1$ if edge $e$ is included in the subgraph, and $x_e = 0$ otherwise. The vector $x = (x_e)_{e \in E}$ represents a candidate edge selection.

\paragraph{QUBO Derivation.}
\textbf{Step 1 -- Objective Function.} The original optimization problem is to minimize the total weight of selected edges:
\begin{align}
\min_{x \in \{0,1\}^{|E|}} \sum_{e \in E} w_e x_e.
\end{align}

\textbf{Step 2 -- Connectivity Constraints.} Connectivity is enforced via the cut formulation. Let $\mathcal{C}$ denote the collection of all valid vertex partitions $S \subset V$ such that $S \cap T \neq \emptyset$ and $(V \setminus S) \cap T \neq \emptyset$. For each $S \in \mathcal{C}$, let $\delta(S)$ denote the set of edges crossing the cut. The requirement that all terminals remain connected is equivalent to:
\begin{align}
\sum_{e \in \delta(S)} x_e \geq 1, \quad \forall S \in \mathcal{C}.
\end{align}

\textbf{Step 3 -- Penalty Term Conversion.} Each inequality constraint is converted to a quadratic penalty term with weight $\lambda > 0$. For a fixed cut $S \in \mathcal{C}$, the penalty is:
\begin{align}
P_S(x) = \lambda \left( 1 - \sum_{e \in \delta(S)} x_e \right)^2.
\end{align}
Expanding the square and applying the idempotent property $x_e^2 = x_e$ for binary variables yields:
\begin{align}
P_S(x) &= \lambda \left( 1 - 2\sum_{e \in \delta(S)} x_e + \sum_{e \in \delta(S)} x_e^2 + \sum_{e \neq f \in \delta(S)} x_e x_f \right) \\
&= \lambda \left( 1 - \sum_{e \in \delta(S)} x_e + \sum_{e \neq f \in \delta(S)} x_e x_f \right).
\end{align}
This expression provides the general algebraic contribution of cut $S$ to the QUBO Hamiltonian.

\textbf{Step 4 -- Combined Hamiltonian.} Summing the objective and all cut penalties gives the total energy function $H(x)$:
\begin{align}
H(x) = \sum_{e \in E} w_e x_e + \sum_{S \in \mathcal{C}} \lambda \left( 1 - \sum_{e \in \delta(S)} x_e + \sum_{e \neq f \in \delta(S)} x_e x_f \right).
\end{align}

\textbf{Step 5 -- Q Matrix and Offset.} Collecting linear and quadratic terms into the standard QUBO form $H(x) = \sum_{e \in E} Q_{e,e} x_e + \sum_{e < f} Q_{e,f} x_e x_f + c$, we obtain the following closed-form expressions for the matrix entries and constant offset:
\begin{align}
Q_{e,e} &= w_e - \lambda \sum_{S \in \mathcal{C}} \mathbb{I}(e \in \delta(S)), & \forall e \in E, \\
Q_{e,f} &= \lambda \sum_{S \in \mathcal{C}} \mathbb{I}(e \in \delta(S) \land f \in \delta(S)), & \forall e, f \in E, e < f, \\
c &= \lambda |\mathcal{C}|. &
\end{align}
Here, $\mathbb{I}(\cdot)$ denotes the indicator function, and $|\mathcal{C}|$ is the total number of valid cuts.

\paragraph{Penalty Weight Justification.}
Let $W_{\max} = \max_{e \in E} w_e$ be the maximum edge weight in the graph. The maximum possible reduction in the objective function achievable by selecting any single edge is bounded by $W_{\max}$. To guarantee that any infeasible solution (violating at least one cut constraint) has a strictly higher energy than the optimal feasible solution, the penalty weight must dominate the objective scale. Specifically, choosing $\lambda > \sum_{e \in E} w_e$ ensures that the penalty incurred from a single violated constraint exceeds the total weight of any feasible Steiner tree, thereby forcing the optimizer toward feasible configurations.

\paragraph{Correctness Argument.}
Minimizing $H(x)$ over $\{0,1\}^{|E|}$ recovers the optimal Steiner tree because (i) any feasible solution satisfies all cut constraints, yielding zero penalty, so its energy exactly equals its objective value; (ii) any infeasible solution violates at least one cut constraint, incurring a penalty of at least $\lambda$. Since $\lambda$ is chosen to strictly exceed the maximum possible objective value of any feasible solution, the energy of every infeasible configuration is strictly greater than that of the optimal feasible configuration. Consequently, the global minimum of $H(x)$ corresponds uniquely to the minimum-cost connected subgraph spanning $T$.

\paragraph{Illustrative Example.}
Consider a concrete instance with $M=3$ edges indexed $0, 1, 2$ and weights $w_0, w_1, w_2$. Suppose the valid cuts reduce to two constraints: $\delta(S_1) = \{e_0, e_2\}$ and $\delta(S_2) = \{e_1, e_2\}$. Instantiating the general formulas from Step 4 yields:
\begin{align}
H(x) &= w_0 x_0 + w_1 x_1 + w_2 x_2 + \lambda \left( 1 - x_0 - x_2 + x_0 x_2 \right) + \lambda \left( 1 - x_1 - x_2 + x_1 x_2 \right).
\end{align}
Collecting terms into the standard QUBO form gives:
\begin{align}
H(x) &= (w_0 - \lambda)x_0 + (w_1 - \lambda)x_1 + (w_2 - 2\lambda)x_2 + \lambda x_0 x_2 + \lambda x_1 x_2 + 2\lambda.
\end{align}
The corresponding Q matrix entries and offset are:
\begin{align}
Q_{0,0} &= w_0 - \lambda, & Q_{1,1} &= w_1 - \lambda, & Q_{2,2} &= w_2 - 2\lambda, \\
Q_{0,2} &= \lambda, & Q_{1,2} &= \lambda, & c &= 2\lambda.
\end{align}
All other $Q_{i,j}$ entries are zero. For $\lambda > \max\{w_0, w_1, w_2\}$, the global minimum of this Hamiltonian over $\{0,1\}^3$ identifies the minimum-cost edge subset that connects the terminals, with energy equal to the optimal objective value plus the constant offset $2\lambda$.

\end{document}